\documentclass[11pt,a4paper]{article}

\usepackage[utf8]{inputenc}
\usepackage[T1]{fontenc}
\usepackage{amsmath,amssymb,amsthm}
\usepackage{mathtools}
\usepackage{booktabs}
\usepackage{array}
\usepackage{hyperref}
\usepackage{cleveref}
\usepackage{enumitem}
\usepackage{geometry}
\usepackage{xcolor}
\usepackage{tcolorbox}
\usepackage{listings}

\geometry{margin=2.5cm}

\newtheorem{theorem}{Theorem}[section]
\newtheorem{lemma}[theorem]{Lemma}
\newtheorem{proposition}[theorem]{Proposition}
\newtheorem{corollary}[theorem]{Corollary}
\theoremstyle{definition}
\newtheorem{definition}[theorem]{Definition}
\newtheorem{example}[theorem]{Example}
\theoremstyle{remark}
\newtheorem{remark}[theorem]{Remark}

\lstdefinestyle{coq}{
  language=,
  basicstyle=\ttfamily\small,
  keywordstyle=\bfseries\color{blue!70!black},
  commentstyle=\itshape\color{green!50!black},
  morekeywords={Definition,Theorem,Lemma,Proof,Qed,Fixpoint,Record,Inductive,
    forall,exists,fun,match,with,end,let,in,if,then,else,From,Require,Import},
  breaklines=true, frame=single, backgroundcolor=\color{gray!5},
  xleftmargin=1em, framexleftmargin=0.5em,
}
\lstset{style=coq}

\newtcolorbox{keyresult}[1][]{
  colback=blue!3, colframe=blue!50!black,
  fonttitle=\bfseries, title={#1}, boxrule=0.5pt, arc=2pt
}

\newcommand{\RP}{\texttt{RealProcess}}
\newcommand{\natQ}{\mathbb{N} \to \mathbb{Q}}
\newcommand{\Qed}{\textsc{Qed}}

% ══════════════════════════════════════════════════════
\begin{document}

\title{%
  \textbf{Process Mathematics:\\
  Classical Analysis Without Completed Infinity}\\[0.5em]
  \large Formal Verification in Rocq with 21,603 Theorems
}

\author{
  Horsocrates\\
  \small Independent Researcher\\
  \small \texttt{github.com/Horsocrates/theory-of-systems-coq}
}

\date{April 2026}

\maketitle

% ══════════════════════════════════════════════════════
\begin{abstract}
We develop a formal framework for classical mathematical analysis in which infinity is treated as a property of processes rather than completed objects.
The central construction is $\RP := \natQ$---a function from natural numbers to rationals---which replaces the real number line~$\mathbb{R}$.

Within this framework, formalized in the Rocq proof assistant (v9.0.1), we establish:
\begin{enumerate}[nosep]
\item Fifteen theorems from Wiedijk's ``100 theorems'' benchmark, plus seven additional results of comparable significance (IVT, EVT, Bolzano--Weierstrass, Heine--Borel, FTC, MVT, Taylor, Picard--Lindel\"of, Stone--Weierstrass, and others), all without the Axiom of Infinity or Axiom of Choice.
\item Process completeness: every Cauchy sequence of processes converges stagewise to a diagonal process, with an explicit metric and convergence proof.  We analyze the relationship to classical completeness and state precisely what the Cauchy-completeness bridge does and does not provide.
\item Proof-theoretic analysis: ordinals to~$\varepsilon_0$ with well-foundedness proved as a theorem, and an explicit discussion of how P4 relates to Church's Thesis and to the strength of the ambient proof assistant.
\end{enumerate}

\noindent
The formalization comprises 21,603 machine-verified theorems across 1,483 files with 0~Admitted and 2~axioms (classical logic and constructive witness extraction).

\medskip\noindent
\textbf{Companion papers:} The physics applications (E/R/R framework, Weinberg angle, Born rule) and the data compression pipeline are treated in separate papers~\cite{physics_paper,compression_paper}.

\medskip\noindent
\textbf{Keywords:} process mathematics, constructive analysis, formal verification, Rocq, Cauchy sequences, P4, ordinals, Wiedijk benchmark
\end{abstract}

\tableofcontents
\newpage

% ══════════════════════════════════════════════════════
\section{Introduction}
\label{sec:intro}
% ══════════════════════════════════════════════════════

\subsection{The Problem}

Standard mathematical analysis rests on the real number line~$\mathbb{R}$, conceived as a completed, uncountable set.  This construction requires the Axiom of Infinity (to assert the existence of~$\mathbb{N}$ as a completed set) and typically the Axiom of Choice.  The Theory of Systems~\cite{tos2025} proposes an alternative: derive mathematics from the principle $A = \text{exists}$, with infinity treated as a property of \emph{processes} (Principle~P4), not of completed objects.

\subsection{Scope of This Paper}

This paper establishes the \emph{mathematical} foundations: the universal type \RP{}, process completeness, proof-theoretic strength, and classical analysis theorems.  We do not treat the physics applications or the data compression pipeline; these are the subject of companion papers~\cite{physics_paper,compression_paper}.

\subsection{Contributions}

\begin{enumerate}
\item \textbf{P4 as ontological position} (\S\ref{sec:type}).  We present P4 transparently: it is equivalent to a finitist ontology augmented with classical logic.  We analyze its relationship to Church's Thesis and state precisely what it entails for the function space $\mathbb{N} \to \mathbb{N}$.

\item \textbf{Process completeness} (\S\ref{sec:completeness}).  We define a metric on \RP{}, construct the diagonal limit, prove stagewise completeness, and give an honest analysis of what this does and does not recover from classical completeness.

\item \textbf{Proof-theoretic calibration} (\S\ref{sec:strength}).  We reach~$\varepsilon_0$, prove well-foundedness, and discuss the relationship between P4's claimed restrictions and Rocq's actual strength.

\item \textbf{Classical analysis without Infinity} (\S\ref{sec:theorems}).  Nine theorems in the unified E/R/R pattern: process + Cauchy + bound = existence.

\item \textbf{Wiedijk benchmark} (\S\ref{sec:wiedijk}).  Fifteen of the 100 theorems, all without the Axiom of Infinity.
\end{enumerate}


% ══════════════════════════════════════════════════════
\section{The Universal Type and P4}
\label{sec:type}
% ══════════════════════════════════════════════════════

\subsection{RealProcess}

\begin{definition}[RealProcess]
A \emph{real process} is a function $R : \mathbb{N} \to \mathbb{Q}$.  We write $\RP := \natQ$.
\end{definition}

\begin{definition}[Cauchy condition]
A real process~$R$ is \emph{Cauchy} if for every $\varepsilon \in \mathbb{Q}$, $\varepsilon > 0$, there exists $N \in \mathbb{N}$ such that for all $m, n \ge N$:
$|R(m) - R(n)| < \varepsilon$.
\end{definition}

In Rocq:
\begin{lstlisting}
Definition RealProcess := nat -> Q.

Definition is_Cauchy (R : RealProcess) : Prop :=
  forall eps : Q, 0 < eps ->
    exists N : nat, forall m n : nat,
      (N <= m)%nat -> (N <= n)%nat ->
      Qabs (R m - R n) < eps.
\end{lstlisting}

Under P4, the Cauchy process \emph{is} the real number.  There is no separate limit object.

\subsection{P4 as Ontological Position}
\label{sec:p4position}

We state P4 precisely and transparently.

\begin{quote}
\textbf{P4 (Finite Actuality):} At any stage, finitely many objects are actual.  The type $\texttt{nat}$ is an inductive type (a process of succession), not a completed set.  Every function $\mathbb{N} \to \mathbb{N}$ that exists is given by a finite description.
\end{quote}

\textbf{Relationship to Church's Thesis.}  The last clause---every function $\mathbb{N} \to \mathbb{N}$ is finitely described---is equivalent to the \emph{ontological} form of Church's Thesis: every total function on natural numbers is computable.  We adopt this as an explicit philosophical position, not as a derived consequence.  The collapse of $\Pi^1_1$ quantification to arithmetic (\S\ref{sec:pi11}) is then a direct corollary, not a surprise.

\textbf{What P4 entails:}
\begin{itemize}[nosep]
\item The Axiom of Infinity (a completed $\mathbb{N}$ exists) is rejected.
\item The full Axiom of Choice (selection from uncountable families) is rejected.
\item Countable choice (selection from a process of finite stages) is compatible.
\item The law of excluded middle (L3) is retained---P4 is not constructivism.
\end{itemize}

\textbf{What P4 does \emph{not} entail by itself.}  P4 alone does not determine the proof-theoretic strength of the system.  That depends on which inductive types and recursion principles are admitted.  See~\S\ref{sec:rocqstrength} for a discussion.

\subsection{Incompatibility with Completed Infinity}

\begin{definition}
$\mathrm{CompletedInfSet}(S) :\Leftrightarrow \forall n : \mathbb{N},\; S(n)$.
\end{definition}

\begin{definition}
$\mathrm{P4\_bounded}(\mathrm{actual}) :\Leftrightarrow \forall \text{stage},\; \exists \text{bound},\; \forall n,\; \mathrm{actual}(\text{stage}, n) \Rightarrow n \le \text{bound}$.
\end{definition}

\begin{theorem}\label{thm:p4prohibit}
If $S$ is a completed infinite set, $\mathrm{actual}$ is P4-bounded, and every member of~$S$ is actual at stage~0, then $\bot$.
\end{theorem}

\begin{proof}
P4 gives bound~$B$ at stage~0.  $S(B{+}1)$ holds.  The bridge makes $B{+}1$ actual.  But $B{+}1 > B$, contradicting the bound.
\end{proof}

\begin{remark}
The force of this theorem depends on the bridge assumption: ``every member of~$S$ is actual at stage~0.''  This is not a logical necessity but an ontological commitment: if a completed set exists, its members are simultaneously actual.  Under P4, this commitment is precisely what is rejected.  The theorem shows that P4-boundedness and the bridge are \emph{jointly} inconsistent with completed infinity---not that completed infinity is ``refuted'' in any absolute sense.
\end{remark}

\emph{Rocq}: \texttt{P4CompletedInfinity.v}, theorem \texttt{completed\_inf\_contradicts\_P4}, \Qed.


% ══════════════════════════════════════════════════════
\section{Process Completeness}
\label{sec:completeness}
% ══════════════════════════════════════════════════════

This section addresses the central objection: ``Without~$\mathbb{R}$, Cauchy sequences may not converge.''  We define a metric, construct limits, and analyze precisely what is and is not recovered.

\subsection{Stage Distance}

\begin{definition}[Stage distance]
For processes $R, S : \RP$ and truncation level~$N$:
\[
  d_N(R, S) = \sum_{k=0}^{N-1} |R(k) - S(k)|.
\]
\end{definition}

Properties verified in Rocq (\texttt{ProcessMetricComplete.v}):
\begin{align}
  d_N(R, R) &= 0, \label{eq:self}\\
  d_N(R, S) &\ge 0, \label{eq:nonneg}\\
  d_N(R, S) &\le d_{N+1}(R, S). \label{eq:mono}
\end{align}

Under P4, the ``full'' distance $d(R, S) = \lim_{N \to \infty} d_N(R, S)$ does not exist as a completed object.  The sequence $\{d_N(R, S)\}$ is itself a process over~$\mathbb{Q}$.

\subsection{Process-Cauchy and Diagonal Limit}

\begin{definition}[Process-Cauchy]
A sequence $\mathrm{seq} : \mathbb{N} \to \RP$ is \emph{process-Cauchy} if for every truncation~$N$ and every $\varepsilon > 0$, there exists~$M$ such that for all $i, j \ge M$:
$d_N(\mathrm{seq}(i), \mathrm{seq}(j)) < \varepsilon$.
\end{definition}

\begin{definition}[Diagonal process]
$R_\infty(k) = \mathrm{seq}(k)(k)$.
\end{definition}

\begin{theorem}[Stagewise completeness]\label{thm:completeness}
If $\mathrm{seq}$ is process-Cauchy, then for every stage~$k$ and $\varepsilon > 0$, there exists~$M$ such that for all $i \ge M$ with $i \ge k$:
\[
  |\mathrm{seq}(i)(k) - R_\infty(k)| < \varepsilon.
\]
\end{theorem}

\begin{proof}
Process-Cauchy at $N = k{+}1$ gives $M_k$ with $|\mathrm{seq}(i)(k) - \mathrm{seq}(j)(k)| < \varepsilon$ for $i, j \ge M_k$.  Take $j = k$ (requiring $k \ge M_k$, so set $M = \max(M_k, k)$):
$|\mathrm{seq}(i)(k) - \mathrm{seq}(k)(k)| = |\mathrm{seq}(i)(k) - R_\infty(k)| < \varepsilon$. \qedhere
\end{proof}

\emph{Rocq}: \texttt{ProcessMetricComplete.v}, theorem \texttt{process\_completeness}, \Qed.

\subsection{Relationship to Classical Completeness}
\label{sec:bridge}

Stagewise completeness says: for each \emph{fixed}~$k$, the $k$-th components converge.  This is \emph{not} the same as convergence in a completed metric space.  We now state precisely what the bridge to classical completeness provides.

\textbf{What \texttt{ProcessBridge.v} establishes:}
\begin{enumerate}[nosep]
\item A Cauchy process $R : \natQ$ can be packaged as a classical Cauchy sequence of rationals (a \texttt{CauchySeq} record in Rocq).
\item If classical $\mathbb{R}$ is available (e.g., via Rocq's \texttt{Reals} library), then every \texttt{CauchySeq} has a limit in~$\mathbb{R}$.
\item The composition gives: every Cauchy process determines a real number.
\end{enumerate}

\textbf{What it does \emph{not} establish:}
\begin{enumerate}[nosep]
\item It does not construct~$\mathbb{R}$ within P4.  Step~2 imports the classical reals.
\item It does not prove that the process metric (stagewise distance) is complete in the classical sense.
\item It does not show that every real number arises from a Cauchy process (only the converse direction is trivial).
\end{enumerate}

\begin{keyresult}[Honest assessment]
Under P4, we have \emph{stagewise completeness}: every Cauchy sequence of processes converges componentwise to a diagonal process.  This suffices for all nine classical analysis theorems in~\S\ref{sec:theorems}: each constructs a Cauchy process and bounds its rate, never needing a completed limit object.

Classical completeness---the existence of a limit \emph{point} in a complete metric space---is available via the bridge, but the bridge imports $\mathbb{R}$ from outside the P4 framework.  Under P4 alone, the process \emph{is} the real number, and stagewise completeness is the operative notion.
\end{keyresult}

\subsection{Consequences for the Function Space}

If every function $\mathbb{N} \to \mathbb{N}$ is computable (the ontological commitment of P4, see~\S\ref{sec:p4position}), then:
\begin{itemize}[nosep]
\item The ``reals'' (Cauchy processes over~$\mathbb{Q}$) are \emph{countable}---only computable Cauchy sequences exist.
\item Continuous-but-nowhere-differentiable functions exist only if computably given.
\item Large parts of descriptive set theory (Borel hierarchy beyond $\Sigma^0_1$, analytic sets) are modified.
\end{itemize}

These are genuine mathematical losses.  We argue (\S\ref{sec:survives}) that they do not affect any theorem used in science or engineering, but we do not claim they are negligible from a mathematical standpoint.


% ══════════════════════════════════════════════════════
\section{Proof-Theoretic Strength}
\label{sec:strength}
% ══════════════════════════════════════════════════════

\subsection{Ordinals to $\varepsilon_0$}

\begin{lstlisting}
Inductive Ord : Set :=
  | OZero : Ord
  | OSucc : Ord -> Ord
  | OLim  : (nat -> Ord) -> Ord.
\end{lstlisting}

Concrete ordinals defined: $\omega = \texttt{OLim}(\texttt{nat\_to\_ord})$, $\omega^\omega$, and $\varepsilon_0 = \texttt{OLim}(\texttt{omega\_tower})$.  Full arithmetic (addition, multiplication, exponentiation) verified.

\subsection{Well-Foundedness as Theorem}

\begin{theorem}\label{thm:wf}
$\texttt{wf\_ord\_lt} : \texttt{well\_founded}\;\texttt{ord\_lt}$.
\end{theorem}

Proved by structural induction on \texttt{Ord} with an accessibility argument (helper \texttt{acc\_succ}), requiring no additional axioms.

\emph{Rocq}: \texttt{TransfiniteInduction.v}, 12~\Qed.

\subsection{ATR$_0$ as Definition}

Under P4, transfinite recursion along \texttt{Ord} is a \texttt{Fixpoint}---accepted by Rocq's termination checker because recursion is structural.  The ``axiom'' ATR$_0$ becomes a function definition.

\emph{Rocq}: \texttt{P4\_Eliminates\_ATR.v}, 15~\Qed.

\subsection{$\Pi^1_1$ and Church's Thesis}
\label{sec:pi11}

Under P4's ontological commitment (every function $\mathbb{N} \to \mathbb{N}$ is finitely described), second-order quantification collapses:
\[
  \forall f : \mathbb{N} \to \mathbb{N}.\; \varphi(f, n)
  \quad\equiv\quad
  \forall c : \mathbb{N}.\; \varphi(\texttt{eval}(c), n).
\]

\textbf{We emphasize:} this is a \emph{corollary} of the P4 ontology (which is equivalent to the ontological form of Church's Thesis), not an independent derivation.  The result is unsurprising given the premise; its value lies in making the consequence explicit.

\emph{Rocq}: \texttt{P4\_Eliminates\_Pi11.v}, 15~\Qed.

\subsection{P4 vs.\ Rocq's Strength}
\label{sec:rocqstrength}

A methodological issue deserves transparency.  Rocq's Calculus of Inductive Constructions (CIC) has proof-theoretic strength far beyond~$\varepsilon_0$---it can define ordinals past~$\Gamma_0$ and beyond.  Our formalization lives \emph{within} Rocq but does not \emph{use} its full strength.

\begin{keyresult}[What we claim and what we don't]
\textbf{We claim:} constructions up to~$\varepsilon_0$ \emph{suffice} for classical analysis (IVT, EVT, BW, FTC, MVT, Taylor, Picard, and 15 Wiedijk theorems).  This is demonstrated by the 21,603~\Qed.

\textbf{We do not claim:} $\varepsilon_0$ is the \emph{maximum} that P4 allows.  A formal characterization of ``P4-admissible inductive types'' within CIC is an open problem.  Our position is: the interesting question is not what P4 \emph{prohibits} but what suffices.
\end{keyresult}

\subsection{Position in the Hierarchy}

\begin{table}[h]
\centering
\begin{tabular}{lccccc}
\toprule
System & L3 & Ax.\ Inf. & AC & Ordinal & Coverage \\
\midrule
PRA         & varies & no  & no  & $\omega^\omega$ & primitive recursive \\
PA          & yes    & no  & no  & $\varepsilon_0$ & first-order arithmetic \\
ACA$_0$     & yes    & yes & no  & $\varepsilon_0$ & predicative analysis \\
ZFC         & yes    & yes & yes & inaccessible & all standard math \\
\midrule
\textbf{ToS} & \textbf{yes} & \textbf{no} & \textbf{no} & $\boldsymbol{\varepsilon_0}$\textsuperscript{*} & \textbf{classical analysis} \\
\bottomrule
\end{tabular}
\caption{Proof-theoretic hierarchy.  \textsuperscript{*}Constructions to~$\varepsilon_0$ demonstrated; upper bound not formally established (see~\S\ref{sec:rocqstrength}).}
\label{tab:hierarchy}
\end{table}


% ══════════════════════════════════════════════════════
\section{Classical Theorems as Process Constructions}
\label{sec:theorems}
% ══════════════════════════════════════════════════════

Every classical existence theorem has the same structure: construct a process over~$\mathbb{Q}$, show it is Cauchy with computable rate, verify the property at every stage.

\subsection{Intermediate Value Theorem}

\begin{theorem}[IVT, Process Version]
Let $f : \mathbb{Q} \to \mathbb{Q}$ be uniformly continuous on $[a, b]$ with $f(a) < 0 < f(b)$.  Then there exists a Cauchy process $c : \natQ$ with $c(n) \in [a, b]$ and $|f(c(n))| \to 0$.
\end{theorem}

\textbf{Construction.}  Bisection: $c(n)$ is the midpoint of a nested interval halving at each step.  Rate: $|c(m) - c(n)| \le (b-a)/2^{\min(m,n)}$.  Uses L3 for the sign test.

\emph{Rocq}: \texttt{IVT\_ERR.v}, 24~\Qed.

\subsection{Extreme Value Theorem}

Grid refinement: $x^*(n) = \arg\max_k f(a + k(b{-}a)/2^n)$ on a $2^n$-point grid.  The supremum process is monotonically non-decreasing.  When multiple points achieve the maximum, L5-Resolution selects the leftmost.

\emph{Rocq}: \texttt{EVT\_ERR.v}, 36~\Qed.

\subsection{Bolzano--Weierstrass}

Nested interval bisection: at each step, at least one half contains infinitely many terms (pigeonhole + L3).  Rate: interval width $(b{-}a)/2^n$.

\emph{Rocq}: \texttt{BolzanoWeierstrass.v}, 27~\Qed.

\subsection{Fundamental Theorem of Calculus}

Riemann sums on $2^n$-point grids form a Cauchy process.  Lipschitz bound: $|F(n{+}1) - F(n)| \le L(b{-}a)^2/2^{n+1}$.

\emph{Rocq}: \texttt{analysis/FTC.v}, 29~\Qed.

\subsection{Further Theorems}

Mean Value Theorem (\texttt{MeanValueTheorem.v}, 18~\Qed), Taylor's Theorem (\texttt{TaylorSeries.v}, 18~\Qed), Picard--Lindel\"of (\texttt{PicardLindelof.v}, 24~\Qed), Heine--Borel (\texttt{HeineBorelComplete.v}, 30~\Qed), Monotone Convergence (\texttt{MonotoneConvergence.v}, 15~\Qed).

\subsection{The Unified Pattern}

\begin{table}[h]
\centering
\small
\begin{tabular}{@{}lllll@{}}
\toprule
Theorem & Elements & Roles & Rules & Rate \\
\midrule
IVT          & midpoints       & Cauchy, bounded  & bisection           & $1/2^n$ \\
EVT          & grid maximizers & monotone         & grid refinement     & $\delta(2^{-n})$ \\
BW           & intervals       & shrinking        & pigeonhole (L3)     & $1/2^n$ \\
Heine--Borel & grid points     & covering         & Lebesgue number     & finite \\
FTC          & Riemann sums    & Cauchy           & Lipschitz bound     & $1/2^n$ \\
MVT          & critical pts    & extremal         & IVT + EVT           & $\delta(2^{-n})$ \\
Taylor       & partial sums    & bounded error    & remainder           & $1/n!$ \\
Picard       & iterates $y_n$  & Cauchy           & contraction         & $(Lt)^n/n!$ \\
Monotone CV  & bounded seq     & converging       & supremum            & $1/n$ \\
\bottomrule
\end{tabular}
\caption{Nine process theorems.  Each: process over~$\mathbb{Q}$, Cauchy condition, explicit rate.}
\label{tab:err}
\end{table}


% ══════════════════════════════════════════════════════
\section{Wiedijk's 100 Theorems}
\label{sec:wiedijk}
% ══════════════════════════════════════════════════════

\begin{table}[h]
\centering
\small
\begin{tabular}{@{}rllr@{}}
\toprule
\# & Theorem & Rocq File & \Qed \\
\midrule
1  & $\sqrt{2}$ irrational & \texttt{Sqrt2Irrational.v} & 14 \\
25 & Schroeder--Bernstein & \texttt{SchroederBernstein\_ERR.v} & 16 \\
34 & Harmonic series diverges & \texttt{HarmonicDiverges.v} & 21 \\
35 & Taylor's theorem & \texttt{TaylorSeries.v} & 18 \\
39 & Fibonacci / golden ratio & 3 files & 34 \\
58 & Binomial coefficients & \texttt{Combinatorics.v} & 26 \\
63 & Cantor's theorem & 2 files & 215 \\
64 & L'H\^opital's rule & \texttt{LHopital.v} & 25 \\
69 & GCD & \texttt{GCD.v} & 23 \\
74 & Cauchy--Schwarz & \texttt{CauchySchwarz.v} & 18 \\
75 & Bolzano--Weierstrass & \texttt{BolzanoWeierstrass.v} & 27 \\
78 & Mean value theorem & \texttt{MeanValueTheorem.v} & 18 \\
79 & Intermediate value thm & \texttt{IVT\_ERR.v} & 24 \\
80 & Fundamental thm calculus & \texttt{FTC.v} & 29 \\
86 & Lebesgue measure & \texttt{LebesgueMeasure.v} & 23 \\
\midrule
\multicolumn{3}{@{}l}{\textbf{Total from Wiedijk's 100:}} & \textbf{531} \\
\bottomrule
\end{tabular}
\caption{Fifteen theorems from Wiedijk's list, all proved over~$\mathbb{Q}$ without the Axiom of Infinity.}
\label{tab:wiedijk}
\end{table}

Seven additional results (EVT, Heine--Borel, Stone--Weierstrass, Picard--Lindel\"of, Implicit Function, Banach Fixed Point, Monotone Convergence) bring the total to 22 theorems and 703~\Qed.


% ══════════════════════════════════════════════════════
\section{What Survives, What Is Lost}
\label{sec:survives}
% ══════════════════════════════════════════════════════

\subsection{What Survives}

All theorems in Tables~\ref{tab:err} and~\ref{tab:wiedijk}, plus the broader formalization (21,603~\Qed{} covering algebra, topology, measure theory, ODEs, spectral theory, and lattice gauge theory).  Every computable mathematical statement that holds in ZFC also holds under P4.

\subsection{What Is Lost---Honestly}

P4 with the Church's Thesis commitment entails:

\begin{enumerate}
\item \textbf{The reals are countable.}  Only computable Cauchy sequences exist.  The classical uncountability of~$\mathbb{R}$ is replaced by the statement: no computable enumeration captures all computable Cauchy sequences (a weaker but still meaningful diagonal result, formalized in \texttt{ShrinkingIntervals\_ERR.v}).

\item \textbf{Non-measurable sets do not exist.}  Every definable set of reals is measurable.  Banach--Tarski is dissolved.  This is widely regarded as a feature, not a bug.

\item \textbf{Descriptive set theory is modified.}  The Borel hierarchy, analytic sets, and the theory of definability change character when only computable objects exist.

\item \textbf{Full Kruskal's tree theorem is unavailable.}  We prove it for specific families (chains, forks, bounded depth) but the full minimal bad sequence argument requires $\Pi^1_2$ strength, which exceeds our demonstrated~$\varepsilon_0$.

\item \textbf{Continuous-but-nowhere-differentiable functions} exist only if given by a computable rule.  The ``generic'' continuous function of classical analysis---which exists by a Baire category argument over non-computable objects---does not exist under P4.
\end{enumerate}

\subsection{Assessment}

No physical prediction, no engineering algorithm, and no computational procedure depends on the existence of non-computable functions or non-measurable sets.  The mathematical losses are real but confined to areas of pure mathematics that concern the behavior of non-constructive objects.  Whether this is an acceptable tradeoff is a philosophical judgment, not a mathematical one.


% ══════════════════════════════════════════════════════
\section{Formalization Statistics}
\label{sec:stats}
% ══════════════════════════════════════════════════════

\begin{table}[h]
\centering
\small
\begin{tabular}{@{}lrr@{}}
\toprule
Domain & Files & \Qed \\
\midrule
Foundation (laws, principles, E/R/R) & 117 & 1,589 \\
Process framework (Phases 0--10) & 297 & 4,180 \\
Standard library & 627 & 7,510 \\
Gauge theory & 114 & 2,177 \\
Classical analysis & 34 & 590 \\
Other (zeta, physics, experimental, \ldots) & 294 & 5,557 \\
\midrule
\textbf{Total} & \textbf{1,483} & \textbf{21,603} \\
\bottomrule
\end{tabular}
\caption{Formalization statistics.  0~Admitted, 2~axioms (\texttt{classic}, \texttt{L4\_witness}).}
\label{tab:stats}
\end{table}

\begin{center}
\small
\begin{tabular}{@{}lrrll@{}}
\toprule
Project & Theorems & Wiedijk & Logic & Ax.\ $\infty$ \\
\midrule
Lean \texttt{mathlib} & 200,000+ & 100 & CIC & yes \\
Isabelle AFP & 150,000+ & 76 & HOL & yes \\
Mizar MML & 60,000+ & 69 & set theory & yes \\
\textbf{ToS} & \textbf{21,603} & \textbf{15} & \textbf{CIC + classic} & \textbf{no} \\
\bottomrule
\end{tabular}
\end{center}

Our contribution is not size but foundations: we demonstrate that classical analysis does not require completed infinities.


% ══════════════════════════════════════════════════════
\section{Conclusion}
\label{sec:conclusion}
% ══════════════════════════════════════════════════════

Classical analysis works without completed infinities.  The type $\RP = \natQ$, combined with a finitist ontology (P4) and classical logic (L3), suffices for 22 classical analysis theorems, 15 from Wiedijk's benchmark, with process completeness via diagonal construction.

The framework occupies a unique position: classical logic without the Axiom of Infinity, reaching~$\varepsilon_0$ in proof-theoretic strength.  Its losses are real (countable reals, no non-measurable sets, modified descriptive set theory) but do not affect computable mathematics.

The companion papers~\cite{physics_paper,compression_paper} demonstrate that the same framework supports physics (Weinberg angle, Born rule, gauge theory) and engineering (300$\times$ IoT compression).


\begin{thebibliography}{99}
\bibitem{tos2025} Horsocrates (2025). Theory of Systems: A First-Principles Foundation for Mathematics. PhilPapers.
\bibitem{rocq} The Rocq Development Team (2025). The Rocq Prover, version 9.0.1.
\bibitem{bishop} Bishop, E. (1967). Foundations of Constructive Analysis. McGraw-Hill.
\bibitem{simpson} Simpson, S. G. (2009). Subsystems of Second Order Arithmetic. CUP.
\bibitem{wiedijk} Wiedijk, F. (2008). Formalizing 100 Theorems. \url{https://www.cs.ru.nl/~freek/100/}
\bibitem{physics_paper} Horsocrates (2026). Three E/R/R Formulas for Physics. In preparation.
\bibitem{compression_paper} Horsocrates (2026). Verified Data Compression from First Principles. In preparation.
\end{thebibliography}

\end{document}
